\chapter[Resumen de Esfuerzo Requerido]{Resumen del Esfuerzo Requerido}

El proyecto planifica 480 horas de trabajo, que resultan de dos desarrolladoras
con 15 horas semanales durante 16 semanas. La Tabla~\ref{tab:esfuerzo} las
distribuye por fase.

\begin{table}[H]
\centering
\small
\tablazebra
\caption{Esfuerzo por fase del desarrollo.}
\label{tab:esfuerzo}
\begin{tabularx}{\textwidth}{P{1.86}P{0.62}P{0.52}}
\toprule
\filaencabezado \encab{Actividad o fase} & \encab{Participación} & \encab{Horas} \\
\midrule
Levantamiento de información con Enfermería & 10~\% & 48 \\
Análisis y especificación de requerimientos & 15~\% & 72 \\
Diseño de la solución y prototipos & 15~\% & 72 \\
Construcción en tres sprints & 40~\% & 192 \\
Pruebas automatizadas & 15~\% & 72 \\
Documentación & 5~\% & 24 \\
\midrule
\filatotal \textbf{Total} & \textbf{100~\%} & \textbf{480} \\
\bottomrule
\end{tabularx}
\end{table}

La estimación por puntos de caso de uso del anexo de planificación arroja
129~UCP. Con el factor de 20~horas por UCP que propone la industria, el esfuerzo
esperado sería de unas 2.580 horas, más de cinco veces el presupuesto disponible.
Esta brecha explica tres decisiones del proyecto. La primera fue acotar el
producto a un prototipo de nivel TRL~3. La segunda fue priorizar el backlog para
construir primero el flujo clínico principal. La tercera fue apoyarse en
paquetes y software libre existentes, como Express, PostgreSQL y las librerías
de cifrado, exportación y notificaciones, que resuelven parte importante de la complejidad.

A la fecha de este informe, la aplicación suma 27.413 líneas de código Dart en
188 archivos y el backend 1.483 líneas de código JavaScript en 16 archivos, más
294 casos de prueba de la aplicación y 54 del backend. La relación efectiva entre el esfuerzo
planificado y el tamaño estimado es de unas 3,7 horas por UCP. Las horas reales
de cada fase se registrarán al cierre del sprint~3 para compararlas con esta
planificación.
